\chapter{Shared Request And Response Conventions}

\section{Versioning And Transport}

All implemented endpoints use the \code{/v1} prefix. Requests and responses are JSON
unless the route is an outbound webhook delivery performed by the gateway
toward a merchant callback URL.

\section{Time, Money, And Identifiers}

\begin{longtable}{L{0.28\textwidth}L{0.62\textwidth}}
\caption{Shared data conventions}\\
\toprule
\textbf{Item} & \textbf{Convention} \\
\midrule
\endfirsthead
\toprule
\textbf{Item} & \textbf{Convention} \\
\midrule
\endhead
Timestamps & UTC ISO 8601 timestamps are used for API payloads and response fields. \\
Amounts & Monetary amounts are represented as decimal strings rather than binary floating-point values. \\
Merchant scope & Merchant-facing HMAC APIs derive merchant identity from auth headers; Merchant Portal APIs derive scope from the authenticated portal session. \\
Audit reason & Internal write routes use a nested actor object so the human reason can be written to audit logs. \\
Scenario IDs & Postman and testing scenarios reuse the same stable IDs documented here. \\
\bottomrule
\end{longtable}

\section{Merchant HMAC Headers}

Merchant-facing payment and refund endpoints require:

\begin{itemize}
  \item \code{X-Merchant-Id}
  \item \code{X-Access-Key}
  \item \code{X-Timestamp}
  \item \code{X-Signature}
  \item \code{X-Idempotency-Key} when technical idempotency is desired
\end{itemize}

Canonical signing string:

\begin{lstlisting}[caption={Merchant HMAC signing string}]
{timestamp}.{method}.{path}.{body_sha256_hex}
\end{lstlisting}

\section{Internal Write Payload Convention}

Many internal Ops write routes use the following nested actor shape so that the
audit reason remains explicit:

\begin{lstlisting}[language=json,caption={Internal actor payload fragment}]
{
  "actor": {
    "actor_type": "OPS",
    "actor_id": null,
    "reason": "Operator note"
  }
}
\end{lstlisting}

The backend treats \code{reason} as meaningful input, but it does not trust
client-supplied identity blindly. Role and user identity are derived from the
authenticated internal session.

\section{Shared Error Envelope}

All API errors use the same envelope shape:

\begin{lstlisting}[language=json,caption={Shared error response shape}]
{
  "error_code": "PAYMENT_ALREADY_SUCCESS",
  "message": "Payment already succeeded for this order.",
  "request_id": "req_...",
  "details": {}
}
\end{lstlisting}

The shared error model is summarized later in this document and corresponds to
\code{docs/api/errors.md}.
